\documentclass[12pt]{article}
\usepackage{fontspec}
\usepackage{polyglossia}
\setmainlanguage{french}
\usepackage{csquotes}

\usepackage{tikz}% faire des schémas
\usepackage{pgfplots}% graphiques à partir de données (courbes, histogrammes)
\pgfplotsset{compat=1.16}
% 5. Sans cette ligne, les années s'écrivent « 1,986 » (séparateur des milliers anglo-saxon).
\pgfplotsset{/pgf/number format/1000 sep=}

\begin{document}

\section{Tracer un diagramme avec pgfplots}

% 2. Chaque point s'écrit (abscisse,ordonnée), à la suite, sans virgule entre les points.
%    Les décimales s'écrivent avec un point : 2.07, et non 2,07.
% 3. Un \addplot par série de données : un pour les hommes, un pour les femmes.
% 4. \legend : une entrée par \addplot, dans le même ordre, séparées par des virgules.
% 5. xtick=data : seules les abscisses des données (1986, 1999, 2010) sont graduées.
% 6. ybar : diagramme en barres verticales (les barres des deux séries se placent côte à côte).
% 7. Le graphique est placé dans un flottant figure, avec une légende.

\begin{figure}[ht]
	\centering
	\begin{tikzpicture}
		\begin{axis}[ybar, xtick=data]
			\addplot coordinates {(1986,2.07) (1999,2.13) (2010,2.13)};
			\addplot coordinates {(1986,5.07) (1999,4.36) (2010,4.01)};
			\legend{Hommes,Femmes}
		\end{axis}
	\end{tikzpicture}
	\caption{Temps de travail domestique quotidien, en heures (Insee)}
\end{figure}

\end{document}
